% A155_Meson_Sektor_De_ch.tex
% Meson-Massen: additiver Ansatz, Geltungsbereich, Sektor-Anschluss
% Autor: Johann Pascher, 23. Juli 2026
% Eigenständiges A-Serie-Dokument. Zwischennummer zwischen A150 und A160.

\section{Zweck und Einordnung}

\textbf{[SETZUNG]} Dieses Dokument behandelt Meson-Massen im
FFGFT-Rahmen mit einem eigenständigen additiven Ansatz. Es klärt dessen
Geltungsbereich, wertet ihn am Pion aus, dokumentiert den Anschluss an
die Z$_3$-Sektor-Struktur (A270) und benennt die strukturelle Grenze,
an der der additive Ansatz endet.

A150 führt den Quarksektor als beschreibende Einordnung, nicht als
Prognosegröße. Dieselbe Vorsicht gilt hier: Quarkmassen sind
schemaabhängig (A150, \S{}Quarks, erster Grund); alle Zahlenwerte
verwenden PDG-Werte ($\overline{\text{MS}}$, 2\,GeV) und tragen deren
Unsicherheit ($\sim$1--2\,\%).

\section{Externe Eingaben}

\textbf{[SETZUNG]} Neben den A-Serie-Größen $\xi$ und
$K_\text{frak} = 1-100\xi$ (Herleitung: A040) verwendet dieses Dokument
drei externe Standardgrößen der QCD; sie werden hier deklariert und
nicht hergeleitet:

\begin{center}
\renewcommand{\arraystretch}{1.3}
{\small\begin{tabular}{llp{6cm}}
\toprule
Größe & Wert & Rolle \\
\midrule
$\Lambda_\text{QCD}$ & $0{,}217$\,GeV & QCD-Skala (PDG-Bereich
$0{,}21$--$0{,}22$\,GeV); einzige dimensionsbehaftete Eingabe des
Bindungsterms \\
$m_u, m_d, m_s$ & $2{,}16$; $4{,}67$; $93{,}4$\,MeV &
Quark-Strommassen ($\overline{\text{MS}}$, 2\,GeV, PDG) \\
Meson-Massen & PDG & Vergleichswerte der Auswertung \\
\bottomrule
\end{tabular}}
\renewcommand{\arraystretch}{1.0}
\end{center}

\section{Der additive Ansatz}

\textbf{[SETZUNG]} Für ein Meson aus zwei Konstituenten wird angesetzt:
\[
  m_M = m_{q_1} + m_{q_2}
  + \Lambda_\text{QCD}\cdot K_\text{frak}^{\,n_\text{eff}}.
\]
Begründung der Form: die Konstituentenmassen gehen additiv ein; der
Bindungsanteil wird als QCD-Skala mal fraktaler Abschwächung
$K_\text{frak}^{\,n_\text{eff}}$ angesetzt --- dieselbe
$K_\text{frak}$-Potenzstruktur, die im Korpus Skalenübergänge trägt
(A040; Ebenen-Übersicht P40, Dok.~190). Der Exponent $n_\text{eff}$
ist zunächst frei und wird am Experiment ausgewertet; seine
topologische Deutung liefert A270.

\section{Geltungsbereich}

\textbf{[K]} Da $K_\text{frak} < 1$, gilt für jedes $n_\text{eff} > 0$:
\[
  \Lambda_\text{QCD}\cdot K_\text{frak}^{\,n_\text{eff}} < \Lambda_\text{QCD}.
\]
\textbf{Der additive Ansatz kann strukturell nur Bindungsanteile
unterhalb $\Lambda_\text{QCD}$ tragen.} Das ist eine Formeigenschaft,
keine Näherungsfrage.

\medskip
\textbf{[K]} Numerische Prüfung:

\begin{center}
\renewcommand{\arraystretch}{1.3}
{\small\begin{tabular}{lcccc l}
\toprule
Meson & $m_M$ (GeV) & $\sum m_q$ & Bindung & $n_\text{eff}$ & Status \\
\midrule
$\pi^0$ ($u\bar u/d\bar d$) & 0{,}1350 & 0{,}0043 & 0{,}1307 & 37{,}79 & im Geltungsbereich \\
$\pi^\pm$ ($u\bar d$)       & 0{,}1396 & 0{,}0068 & 0{,}1327 & 36{,}62 & im Geltungsbereich \\
$K^\pm$ ($u\bar s$)         & 0{,}4937 & 0{,}0956 & 0{,}3981 & --- & Bindung $>\Lambda$ \\
$\eta$                      & 0{,}5479 & 0{,}1868 & 0{,}3611 & --- & Bindung $>\Lambda$ \\
$\rho$ ($u\bar u$)          & 0{,}7753 & 0{,}0043 & 0{,}7709 & --- & Bindung $>\Lambda$ \\
$\omega$ ($u\bar u$)        & 0{,}7827 & 0{,}0043 & 0{,}7783 & --- & Bindung $>\Lambda$ \\
\bottomrule
\end{tabular}}
\renewcommand{\arraystretch}{1.0}
\end{center}

Nur die Pionen liegen im Geltungsbereich. Das Scheitern bei $K$, $\eta$,
$\rho$, $\omega$ ist kein numerischer Defekt, sondern die deklarierte
Grenze der Funktionsform.

\section{Sektor-Anschluss (A270)}

\textbf{[K]} A270 sagt aus der Z$_3$-Sektor-Struktur für Mesonen
(1~getwisteter Sektor) den Exponenten $36+1 = 37$ voraus; der exakte
Referenzwert der Bulk-Relation ist $n = 35{,}99$
($K_\text{frak}^{-36} \approx 16/\pi^2$, A270 \S{}Schlüsselrelation).

Die Auswertung der Pion-Daten liefert:
\[
  n_\text{eff}(\pi^0) = 37{,}79, \qquad n_\text{eff}(\pi^\pm) = 36{,}62.
\]
Beide Werte liegen im Fenster der Sektor-Vorhersage $36+1$.
\textbf{Die Meson-Stufe der Sektor-Tabelle aus A270 ist damit am Pion
empirisch verankert} --- der Exponent wurde hier unabhängig von der
Sektor-Überlegung aus den Massendaten bestimmt.

\section{GMOR-Prüfung und die Grenze des additiven Ansatzes}

\textbf{[K]} Für Pseudo-Goldstone-Bosonen gilt die
Gell-Mann--Oakes--Renner-Relation der chiralen Störungstheorie
(externe Standardphysik):
\[
  m_P^2 = (m_{q_1} + m_{q_2})\cdot B,
  \qquad B = \frac{-2\langle\bar qq\rangle}{f_\pi^2}.
\]
Auswertung mit den externen Eingaben:
\[
  B_\pi = \frac{m_{\pi^\pm}^2}{m_u+m_d} = 2{,}852\ \text{GeV},
  \qquad
  B_K = \frac{m_K^2}{m_u+m_s} = 2{,}550\ \text{GeV},
  \qquad
  \frac{B_\pi}{B_K} = 1{,}12.
\]
Pion und Kaon teilen dieselbe chirale Struktur ($B$ bis auf die bekannte
SU(3)-Brechung konstant). \textbf{Der additive Ansatz versagt beim Kaon
also nicht wegen anderer Physik, sondern wegen falscher Funktionsform}:
linear statt quadratisch (GMOR).

\medskip
\textbf{[S]} \textbf{Kandidat: $B$ trägt den Sektor-Faktor.}
Numerisch fällt auf:
\[
  \frac{B_\pi}{\Lambda_\text{QCD}\cdot K_\text{frak}^{-37}}
  = 7{,}999 \approx 8,
  \qquad\text{d.h.}\qquad
  B \approx 8\,\Lambda_\text{QCD}\cdot K_\text{frak}^{-37}.
\]
Die~8 hat zwei mögliche Korpus-Lesarten: $\dim\text{SU(3)} = 8$
(Anzahl Gluonen --- das chirale Kondensat ist gluonisch getragen) oder
das D4-Verhältnis Weyl-Ordnung/Kissing $= 192/24 = 8$ (A270).
\textbf{P35 greift}: die zentrale Abweichung ($0{,}02\,\%$) liegt weit
unter der Schema-Unsicherheit der Quarkmassen ($\sim$2\,\%); der
Treffer ist innerhalb der Unsicherheit, aber die Präzision ist nicht
belastbar. Gebucht als Kandidat, nicht als Ergebnis.

Wäre der Kandidat tragfähig, ergäbe sich eine GMOR-Form mit
Sektor-Faktor,
\[
  m_P^2 = (m_{q_1}+m_{q_2})\cdot 8\,\Lambda_\text{QCD}
  \cdot K_\text{frak}^{-37},
\]
die Pion \emph{und} Kaon abdeckt (Kaon dann auf $\sim$6\,\% via
SU(3)-Brechung in $B$). Die Prüfung dieser Form auf [K]-Niveau setzt
schema-fixierte Quarkmassen voraus --- offen.

\medskip
\textbf{[SETZUNG]} \textbf{Vektormesonen} ($\rho$, $\omega$) sind keine
Goldstone-Bosonen; weder der additive Ansatz noch GMOR ist für sie
zuständig. Ihre Massen liegen an der chiralen Symmetriebrechungsskala
selbst ($\sim 4\pi f_\pi$). Eine FFGFT-Behandlung des Vektorsektors
existiert nicht und wird hier nicht behauptet.

\section{Baryon-Kandidat: vollständig geometrische Proton-Formel}

\textbf{[S]} Der additive Meson-Ansatz und die Sektor-Struktur (A270)
legen die Frage nahe, ob auch die Proton-Masse eine geometrische
Zerlegung besitzt. Kandidat:
\[
  m_p = \underbrace{\frac{\pi^2}{2}}_{\text{Vol}(B^4)}
  \cdot \underbrace{\frac{\pi}{6}}_{\eta_\text{sc}}
  \cdot \underbrace{N_c\,(1+\alpha_s)}_{\text{Farbe/Kopplung}}
  \cdot \underbrace{e^{-3\xi/4}}_{\xi\text{-Dämpfung}}
  \cdot \underbrace{\frac{\Lambda_\text{QCD}}{2}}_{\text{Skala}}
  = 0{,}9402\ \text{GeV}
\]
gegen PDG $0{,}93827$ ($\Delta = +0{,}21\,\%$). Jeder Faktor ist
gedeutet:

\begin{center}
\renewcommand{\arraystretch}{1.3}
{\small\begin{tabular}{llp{6.2cm}}
\toprule
Faktor & Wert & Deutung \\
\midrule
$\pi^2/2$ & $4{,}9348$ & Volumen der Einheits-4-Kugel (T$^4$-nativ) \\
$\pi/6$ & $0{,}5236$ & Packungsdichte des einfach-kubischen Gitters,
$\text{Vol}(B^3)/2^3$ --- 3D-Analogon zur D4-Dichte $\pi^2/16$ (4D),
die den Bulk-Exponenten trägt (A270) \\
$N_c(1+\alpha_s)$ & $3{,}354$ & Farbzahl mal Kopplungskorrektur
(externe QCD-Größen, \S{}Externe Eingaben) \\
$e^{-3\xi/4}$ & $0{,}9999$ & $\xi$-Dämpfung ($N_c$-fach) \\
$\Lambda_\text{QCD}/2$ & $0{,}1085$\,GeV & halbierte QCD-Skala;
einziger Einheitenträger \\
\bottomrule
\end{tabular}}
\renewcommand{\arraystretch}{1.0}
\end{center}

\textbf{Absorbierbarkeit der Restabweichung:} $+0{,}21\,\%$ schließt
exakt mit $\alpha_s = 0{,}1156$ (statt $0{,}118$; $-2\,\%$) oder
$\Lambda_\text{QCD} = 0{,}2165$ (statt $0{,}217$; $-0{,}2\,\%$) ---
beides tief innerhalb der Unsicherheit der externen Eingaben.

\medskip
\textbf{[S]} \textbf{Strukturverbindung zur Higgs-Formel.} Die beiden
Geometriefaktoren verbinden sich zu
\[
  \frac{\pi^2}{2}\cdot\frac{\pi}{6} = \frac{\pi^3}{12},
  \qquad\text{und}\qquad
  \frac{16\pi^3}{\pi^3/12} = 192
\]
--- die D4-Weyl-Ordnung. Die Proton-Geometrie und die Higgs-Formel
$16\pi^3$ (Dok.~190, K1) stehen damit im exakten Verhältnis der
Weyl-Zahl 192. Nach P35 notiert, nicht gebucht.

\medskip
\textbf{P35-Vorbehalt (konkurrierender Kandidat):} Das durch die
Rechnung fixierte Produkt der freien Faktoren beträgt $X = 0{,}52248$;
$\pi/6 = 0{,}52360$ trifft mit $+0{,}21\,\%$. Numerisch näher läge
$K_\text{frak}^{48{,}4}$ ($-0{,}05\,\%$), jedoch ohne ganzzahligen
Exponenten, ohne Deutung und ohne Absorbierbarkeits-Argument
($K_\text{frak}^{48}$ verfehlt mit $+0{,}49\,\%$). Zwei Kandidaten,
keine Vorwärts-Herleitung: [S], nicht [K]. Der historische Anlass
dieser Analyse (Korrektur einer Altkorpus-Rechnung) ist in Dok.~190,
R61 registriert.

\section{Schichtstatus-Zusammenfassung}

\begin{center}
\renewcommand{\arraystretch}{1.3}
{\small\begin{tabular}{ll}
\toprule
Aussage & Status \\
\midrule
Additiver Ansatz $m_M = \sum m_q + \Lambda\,K_\text{frak}^{n}$ & [SETZUNG] (Ansatz) \\
Externe Eingaben $\Lambda_\text{QCD}$, Quarkmassen & [SETZUNG] (deklariert) \\
Geltungsbereich: nur Bindung $<\Lambda_\text{QCD}$ & [K] (Formeigenschaft) \\
Pion im Geltungsbereich, $n_\text{eff}(\pi^0)=37{,}79$ & [K] \\
Sektor-Anschluss $n_\text{eff}\approx 37 = 36+1$ (A270) & [K] \\
GMOR-Konsistenz $B_\pi \approx B_K$ & [K] (Standard-ChPT) \\
$B \approx 8\,\Lambda\cdot K_\text{frak}^{-37}$ & [S] (Kandidat, P35) \\
Vektormesonen außerhalb & [SETZUNG] \\
Proton-Formel $\frac{\pi^3}{12}N_c(1+\alpha_s)e^{-3\xi/4}\frac{\Lambda}{2}$ & [S] (Kandidat, P35) \\
$16\pi^3/(\pi^3/12)=192$ (Higgs-Verbindung) & [S] (notiert, P35) \\
\bottomrule
\end{tabular}}
\renewcommand{\arraystretch}{1.0}
\end{center}

\section{Offene Punkte}

\begin{enumerate}
\item Schema-fixierte Quarkmassen für eine [K]-Prüfung des
  $8\Lambda K_\text{frak}^{-37}$-Kandidaten.
\item Herkunft der~8: Gluonen-Zahl oder D4-Verhältnis --- beide Lesarten
  unverbunden; nach P35 ungebucht, solange keine Vorwärts-Herleitung
  vorliegt.
\item $\pi^\pm$/$\pi^0$-Differenz ($36{,}62$ vs.\ $37{,}79$): trägt die
  elektromagnetische Selbstenergie des geladenen Pions einen eigenen
  Sektor-Beitrag?
\item Vektormeson-Sektor: vollständig offen.
\item Vorwärts-Herleitung des Baryon-Kandidaten: warum $\pi/6$
  (3D-Packung) neben $\text{Vol}(B^4)$ (4D-Volumen)? Entscheidung
  zwischen den Kandidaten $\pi/6$ und $K_\text{frak}$-Potenz steht aus.
\item Zur Korrektur der historischen Altkorpus-Baryon-Rechnung siehe
  Dok.~190, Register R61.
\end{enumerate}

\section{Zeichenerklärung}

Symbole, die in der gesamten A-Serie verwendet werden, sind in A010 erklärt.
Spezifisch für dieses Dokument:

\begin{center}
\renewcommand{\arraystretch}{1.3}
{\small\begin{tabular}{lp{9cm}}
\toprule
Symbol & Bedeutung \\
\midrule
$m_M$, $m_P$ & Meson-Masse allgemein bzw.\ pseudoskalar \\
$m_{q_i}$ & Quark-Strommassen ($\overline{\text{MS}}$, 2\,GeV, PDG; \S{}Externe Eingaben) \\
$\Lambda_\text{QCD}=0{,}217$\,GeV & QCD-Skala (\S{}Externe Eingaben) \\
$n_\text{eff}$ & Exponent der fraktalen Bindungskorrektur $K_\text{frak}^{\,n_\text{eff}}$ \\
$B$ & GMOR-Parameter, $B=-2\langle\bar qq\rangle/f_\pi^2$; aus $m_P^2=(m_{q_1}+m_{q_2})\,B$ \\
$\langle\bar qq\rangle$, $f_\pi$ & chirales Kondensat; Pion-Zerfallskonstante \\
$\eta_\text{sc}=\pi/6$ & Packungsdichte des einfach-kubischen Gitters \\
$\text{Vol}(B^4)=\pi^2/2$ & Volumen der Einheits-4-Kugel \\
$N_c=3$, $\alpha_s=0{,}118$ & Farbzahl; starke Kopplung (externe QCD-Größen) \\
\bottomrule
\end{tabular}}
\renewcommand{\arraystretch}{1.0}
\end{center}

\medskip\noindent\textbf{Prüfskript:} \texttt{python/A\_Serie\_Skripte/a155\_meson.py}
--- alle numerischen Aussagen dieses Dokuments werden dort mit Assertions geprüft.

\section{Bezüge}

\begin{center}
{\small\begin{tabular}{ll}
\toprule
Dokument & Bezug \\
\midrule
A010 & Zeichen der A-Serie \\
A040 & Herleitung $K_\text{frak} = 1-100\xi$ \\
A110 & Zirkulant (Z$_3$-Algebra des ungetwisteten Sektors) \\
A150 & Quarksektor-Vorsicht: Schemaabhängigkeit, keine Prognosegröße \\
A270 & Z$_3$-Sektor-Struktur, Meson-Exponent $36+1=37$ \\
Dok.~190 & P35, P40; Register R61 (Baryon-Korrektur, Altkorpus) \\
\bottomrule
\end{tabular}}
\end{center}
